\documentclass[10pt,a4paper]{amsart}
\usepackage[margin=19mm]{geometry}
\usepackage{amsmath,amssymb,mathtools}
\usepackage[scr=rsfs,cal=euler]{mathalfa}
\usepackage{xfrac}
\usepackage[unicode,hidelinks,bookmarks=false]{hyperref}
\newcommand{\ie}{i.e.\ }
\newcommand{\cf}{cf.\ }
\newcommand{\resp}{respectively\ }
\newcommand{\Subsection}{\subsection}
\providecommand{\nobreakdash}{\nobreak\hbox{-}}
\setlength{\emergencystretch}{2em}
\multlinegap0pt
\usepackage{bm}
\usepackage{bbm}
\usepackage{dsfont}
\usepackage{xspace}
\usepackage{cancel}
\usepackage{mathtools}
\usepackage{amsthm}
\makeatletter
\@ifundefined{theorem}{\newtheorem{theorem}{Theorem}}{}
\makeatother
\DeclareMathOperator{\diag}{\mathrm{diag}}
\title[Dynamical systems on torus related to general Heun equations]{Dynamical systems on torus related to general Heun equations: phase-lock areas and constriction breaking}
\author{Artem Alexandrov, Alexey Glutsyuk}
\date{Source: arXiv:2507.07282v3 (CC BY 4.0)}
\begin{document}
\maketitle

\noindent\emph{Source and licence.} Dynamical systems on torus related to general Heun equations: phase-lock areas and constriction breaking. Version arXiv:2507.07282v3 (CC BY 4.0), distributed under the Creative Commons Attribution 4.0 International licence, as recorded in the arXiv licence metadata for this exact version and shown on the abstract page https://arxiv.org/abs/2507.07282v3. Full rights notice: licenses/arxiv-2507.07282-CC-BY-4.0.txt.\par\smallskip

\noindent\textit{Editorial scope.} Counted unit: the Theorem whose source label is \texttt{growd3} in the article (source0.tex lines 774--777, source label \texttt{growd3}), delivered together with the article's own complete original proof of it (source0.tex lines 779--842, one \texttt{proof} environment of 64 lines). No mathematics is written, repaired or re-derived: statement and proof are the article's own text line for line. Required context retained uncounted: dyn3 (lines 193--195); montriv (lines 717--717). Every definition, display and label of those lines is retained unchanged. Omitted from this package and not redistributed: 1--192, 196--716, 718--773, 778--778, 843--868. Nothing inside a retained excerpt is omitted or altered: every retained body is delivered exactly as the article's lines are written, so no omission receipt of any kind accompanies this package. Citations: the retained excerpts carry no citation command at all, so no bibliography entry of the article is transcribed and no reference is redistributed. Reference adaptation: none was needed and none was made; every reference inside the delivered unit resolves inside it, so no unresolved reference or citation is produced. Printed slips kept literal: no printed word, symbol, hypothesis or number of the counted statement or of its proof was altered. Layout and class: the article's own class is \texttt{elsarticle} with packages including geometry, amsmath, amssymb, amsthm, mathtools, euscript, mathrsfs, graphicx, wrapfig, xcolor, wasysym, verbatim, empheq, adjustbox; the wrapper is the lane's pinned \texttt{amsart} editorial wrapper at 10pt on A4, which restates the theorem-like environments the retained text needs and adds the article's own macro definitions that the retained text needs (\texttt{\textbackslash diag}), with extra wrapper packages (bm, bbm, dsfont, xspace, cancel, mathtools). The delivered numbering is the unit's own. Assets: no retained span contains a figure environment, a graphics-inclusion command, a PDF-page inclusion, a table or a TikZ picture; no image file is copied into this package, the package declares no asset dependency and no image file is redistributed with it. Licence: version arXiv:2507.07282v3 of this article is distributed under the Creative Commons Attribution 4.0 International licence, as recorded in the arXiv licence metadata for this exact version and shown on the abstract page https://arxiv.org/abs/2507.07282v3. A full rights notice is delivered at licenses/arxiv-2507.07282-CC-BY-4.0.txt. Characters: every retained excerpt is pure ASCII, so the delivered engine, XeLaTeX with its default Latin Modern text font, reports no missing character.
\begin{equation}\label{dyn3}
    \frac{d\theta}{d\tau}=\frac{\cos\theta+B+A\sin\tau}{\omega(1-\delta\cos\tau)}+D, \ \  \omega\in\mathbb R\setminus\{0\}, \ \delta\in(0,1).
\end{equation}

\begin{equation}B^2-1=\omega^2(1-\delta^2)(D-n)^2, \ \ n\in\mathbb{Z}.\label{montriv}\end{equation}

\begin{theorem} \label{growd3}
Consider the family of systems~\eqref{dyn3} with $A=D=0$ $B>1$. Then the monodromy $M$ of the corresponding linear system is trivial if and only if 
$$B^2=\omega^2(1-\delta^2)n^2+1, \ \ \ n\in\mathbb{Z}.$$
\end{theorem}

\begin{proof}

Consider the system with $A=0$ and $D=0$,
\begin{equation}\label{eq:zero-AD-tor-sys}
    \frac{d\theta}{d\tau} = \frac{B+\cos\theta}{\omega(1-\delta\cos\tau)}.
\end{equation}
Let $z=e^{i\tau}$, $\Phi=e^{i\theta}$, which gives us
\begin{equation}\label{eq:Riccati-sys}
    \frac{d\Phi}{dz}=-\frac{1+2B\Phi+\Phi^2}{(\delta(1+z^2)-2z)\omega}.
\end{equation}
It is the Riccati equation obtained by  projectivization of the linear system 
\begin{equation}\label{eq:lin-sys}
    \frac{d}{dz}
    \begin{pmatrix}
        u \\ v
    \end{pmatrix}=
    \frac{1}{(\delta(1+z^2)-2z)\omega}
    \begin{pmatrix}
        2B & 1 \\ -1 & 0
    \end{pmatrix}
    \begin{pmatrix}
        u \\ v
    \end{pmatrix}:
\end{equation}
$\Phi(z)$ is a solution of~\eqref{eq:Riccati-sys}, if and only if $\Phi(z)=\frac{v(z)}{u(z)}$, where $(u(z),v(z))$ is a solution of system~\eqref{eq:lin-sys}. 
We are interested in the monodromy of this linear system, computed on the counterclockwise unit circle contour. To compute it, we need to find poles of the denominator in the right hand side of eq.~\eqref{eq:Riccati-sys}. They are given by
\begin{equation}
    z_{\pm} = \frac{1\pm\sqrt{1-\delta^2}}{\delta}.
\end{equation}
Only the point $z_{-}$ belongs to the unit disc, the integration contour interior. Thus, the point $z_+$ is not relevant for our goals. Denote the matrix in the right-hand side of~\eqref{eq:lin-sys} by $Q$,
\begin{equation}
    Q=
    \begin{pmatrix}
        2B & 1 \\ -1 & 0
    \end{pmatrix}.
\end{equation}
This matrix has the following eigenvalues,
\begin{equation}
    \lambda_{1,2}=B\mp\sqrt{B^2-1}.
\end{equation}
Thus, it is diagonalizable: $\lambda_1\neq \lambda_2$, since $B>1$. 
Since the point $z_{-}$ is a Fuchsian singularity, the monodromy matrix is conjugated to 
\begin{equation}
    M = \exp\left\{2\pi i\mathop{\text{Res}}_{z=z_{-}}\frac{\diag(\lambda_1,\,\lambda_2)}{\omega(\delta(1+z^2)-2z)}\right\}=
    \begin{pmatrix}e^{2\pi i\mu_1} & \\ 0 & e^{2\pi i\mu_2}\end{pmatrix},
\end{equation}
where the quantities $\mu_1$ and $\mu_2$ are defined by
\begin{equation}
    \mu_1=\frac{-B+\sqrt{B^2-1}}{2\omega\sqrt{1-\delta^2}},\quad \mu_2=\frac{-B-\sqrt{B^2-1}}{2\omega\sqrt{1-\delta^2}}.
\end{equation}
The monodromy $M$ is scalar, if and only if the ratio of eigenvalues is equal to unity,
\begin{equation*}
    \frac{e^{2\pi i\mu_1}}{e^{2\pi i\mu_2}}=e^{2\pi i(\mu_1-\mu_2)}=1.
\end{equation*}
This is equivalent to the equality
\begin{equation}
   \mu_1-\mu_2= \frac{\sqrt{B^2-1}}{\omega\sqrt{1-\delta^2}}=n,\quad n\in\mathbb{Z}.
\end{equation}
Rearranging this expression, we find that the monodromy $M$ is trivial,  if and only if the parameters obey
\begin{equation}
    B^2=\omega^2n^2\left(1-\delta^2\right)+1,
\end{equation}
which in fact coincides with eq.~\eqref{montriv} with $D=0$.
\end{proof}

\end{document}
